\documentclass[journal,10pt]{IEEEtran}

\usepackage[utf8]{inputenc}
\usepackage[T1]{fontenc}
\usepackage{amsmath,amssymb}
\usepackage{graphicx}
\usepackage{booktabs}
\usepackage{hyperref}
\usepackage{xcolor}
\usepackage{cite}
\usepackage{multirow}
\usepackage{array}
\usepackage{enumitem}
\usepackage{balance}
\usepackage{url}
\usepackage{pifont}

\newcommand{\ie}{\textit{i.e.}}
\newcommand{\eg}{\textit{e.g.}}
\newcommand{\etal}{\textit{et al.}}
\newcommand{\cmark}{\ding{51}}
\newcommand{\xmark}{\ding{55}}
\newcommand{\TODO}[1]{\textcolor{red}{\textbf{[TODO: #1]}}}
\newcommand{\CITE}{\textcolor{blue}{[?]}}

\begin{document}

\title{A Comprehensive Survey on Large Language Model Quantization:\\
Taxonomy, Techniques, Systems, and Frontiers}

\author{
  Wei~Wang
  \thanks{W. Wang is with the School of Information Engineering, Minzu University of China, Beijing 100081, China, and also affiliated with City University of Macau.}
}

\markboth{IEEE Transactions on XXXX, Vol.~XX, No.~X, 2025}
{Wang: A Comprehensive Survey on Large Language Model Quantization}

\maketitle

\begin{abstract}
Large language models (LLMs) have achieved remarkable progress in language understanding, generation, coding, reasoning, and multimodal interaction. However, their rapidly increasing parameter counts, long-context memory requirements, and autoregressive decoding costs create substantial barriers to practical deployment. Quantization has become one of the most important techniques for reducing memory footprint, inference latency, energy consumption, and serving cost.

Unlike general neural network quantization, LLM quantization faces distinctive challenges, including activation outliers, sensitivity of emergent abilities, long-context key-value (KV) cache growth, prefill--decode phase asymmetry, and the need to preserve instruction-following and safety alignment. This survey presents a taxonomy-centered review of LLM quantization. We organize existing work along four complementary axes: \emph{by stage}, \emph{by target}, \emph{by technique}, and \emph{by scenario}. We further discuss systems, hardware, evaluation protocols, safety, theoretical understanding, and open problems.
\end{abstract}

\begin{IEEEkeywords}
Large language models, model quantization, post-training quantization, quantization-aware training, KV cache quantization, low-precision training, efficient inference.
\end{IEEEkeywords}

\section{Introduction}
\IEEEPARstart{L}{arge} language models have become a central substrate for language understanding, code generation, reasoning, tool use, and multimodal interaction. Their capabilities, however, are accompanied by rapidly increasing computational and memory costs. During training, large models require high-bandwidth accelerators, optimizer states, activation checkpointing, and distributed communication. During inference, the cost is shaped not only by parameter size but also by autoregressive decoding, long-context key-value (KV) cache growth, and the mismatch between compute-bound prefill and memory-bandwidth-bound decode phases. These costs make it difficult to deploy LLMs on edge devices, serve them at high throughput, or scale long-context applications economically.

Quantization is one of the most effective techniques for reducing these costs. It replaces high-precision tensors with lower-precision representations such as INT8, INT4, FP8, FP4, NF4, binary, or ternary formats. In LLMs, quantization can target weights, activations, KV cache tensors, gradients, optimizer states, or communication buffers. It can be applied after training, during fine-tuning, throughout quantization-aware training, or even from the beginning of low-precision pre-training. This flexibility makes quantization a broad design space rather than a single compression trick.

LLM quantization differs substantially from conventional neural network quantization. Decoder-only Transformers exhibit channel-wise activation outliers, heavy-tailed weight distributions, residual-stream sensitivity, and token-dependent activation ranges. Autoregressive decoding further amplifies memory movement because each generated token repeatedly reads model weights and KV cache tensors. Moreover, instruction-following, mathematical reasoning, coding, long-context retrieval, and safety alignment can be more fragile than aggregate perplexity suggests. As a result, a quantization method that appears accurate under short-context language modeling loss may still fail in deployment-oriented workloads.

This survey organizes LLM quantization around four complementary axes. The \emph{stage} axis asks when quantization is introduced: post-training quantization (PTQ), quantization-aware training (QAT), quantized fine-tuning, or low-precision pre-training. The \emph{target} axis asks what tensors are quantized: weights, activations, KV cache, attention states, embeddings, gradients, or optimizer states. The \emph{technique} axis asks how quantization error is controlled: outlier handling, Hessian-guided reconstruction, smoothing, rotation, mixed precision, codebook learning, or data-free calibration. The \emph{scenario} axis asks where the method is used: long-context inference, mixture-of-experts (MoE) models, reasoning models, multimodal LLMs, edge deployment, or high-throughput serving.

\subsection{Scope}
This survey focuses on quantization methods and systems that directly target LLMs or closely related Transformer-based generative models. We cover PTQ, QAT, quantized fine-tuning, low-precision training, weight/activation/KV-cache quantization, outlier mitigation, mixed precision, vector and codebook quantization, serving systems, hardware support, evaluation methodology, safety, and theoretical analysis.

We intentionally exclude domains where quantization has a different primary structure, such as CNN-only quantization, ViT-specific quantization, diffusion model quantization, GNN quantization, robotics/VLA quantization, 3D point cloud quantization, and medical imaging quantization, except when their techniques directly inform LLM quantization.

\subsection{Contributions}
\begin{itemize}[leftmargin=*,nosep]
  \item We present a taxonomy-centered organization of LLM quantization along four axes: stage, target, technique, and scenario.
  \item We connect representative methods such as LLM.int8(), GPTQ, AWQ, SmoothQuant, OmniQuant, QuIP/QuIP\#, SpQR, SqueezeLLM, AQLM, KIVI, KVQuant, QLoRA, QA-LoRA, FP8-LM, and BitNet b1.58 to the specific challenges they address.
  \item We relate algorithmic choices to systems constraints, including prefill--decode asymmetry, quantized kernels, KV cache management, serving frameworks, and hardware-native FP8/FP4 formats.
  \item We summarize evaluation, robustness, safety, and theoretical questions that remain underexplored in current LLM quantization research.
\end{itemize}

\section{Background: LLMs and Quantization Fundamentals}

\subsection{LLM Architecture and Inference}
Most contemporary LLMs are decoder-only Transformers composed of attention blocks, feed-forward networks, normalization layers, residual connections, and token embeddings. Inference consists of two phases with different bottlenecks. The \emph{prefill} phase processes all prompt tokens in parallel and is often compute-bound for long prompts. The \emph{decode} phase generates one token at a time, repeatedly loading model weights and KV cache tensors, and is frequently memory-bandwidth-bound. This phase distinction is central to quantization: weight quantization mainly reduces parameter bandwidth, activation quantization enables low-precision compute, and KV cache quantization reduces memory growth during long-context decoding.

\subsection{Number Formats}
LLM quantization uses both integer and floating-point low-precision formats. INT8 is often used for weight-activation inference when hardware support is mature. INT4 and lower-bit integer formats are common for weight-only compression, especially W4A16 deployment. NF4, popularized by QLoRA~\cite{dettmers2023qlora}, is designed for normally distributed weights in quantized fine-tuning. FP8 formats such as E4M3 and E5M2 are attractive for training and hardware-native inference, while FP4, MXFP4, NVFP4, binary, and ternary formats represent more aggressive frontiers. The choice of format determines not only model quality but also kernel availability, packing overhead, memory alignment, and accelerator support.

\subsection{Quantization Granularity}
Granularity specifies how scales and zero-points are shared. Per-tensor quantization is simple but often too coarse for LLMs. Per-channel and per-group quantization are widely used for weights because they balance accuracy and metadata overhead. Per-token or dynamic activation quantization can better handle input-dependent activation ranges. KV cache quantization may require different granularities for keys and values because their distributions differ across heads, channels, and token positions. Thus, granularity is a core design variable rather than a minor implementation detail.

\subsection{LLM-Specific Challenges}
LLM quantization is dominated by several recurring difficulties. First, activation outliers can force large quantization scales and waste most low-bit levels on rare values. Second, error propagation across residual streams and attention layers can make local quantization loss an unreliable predictor of downstream performance. Third, KV cache memory grows linearly with context length and batch size, making long-context serving a distinct target. Fourth, reasoning, coding, instruction-following, and safety behavior may degrade even when perplexity changes little. Finally, MoE, multimodal, and tool-using LLMs create heterogeneous precision requirements across experts, modalities, and generation phases.

\section{Taxonomy of LLM Quantization}

A useful taxonomy of LLM quantization should answer four questions: \emph{when} quantization is introduced, \emph{what} tensors are quantized, \emph{how} quantization error is controlled, and \emph{where} the resulting model is deployed. These questions define the four axes used throughout this survey.

\begin{itemize}[leftmargin=*,nosep]
  \item \textbf{By Stage}: PTQ, QAT, quantized fine-tuning, and low-precision pre-training. This axis captures the amount of training data, compute, and model access required by a method.
  \item \textbf{By Target}: weights, activations, KV cache, attention states, embeddings, optimizer states, and gradients. This axis captures which memory or compute bottleneck the method attempts to reduce.
  \item \textbf{By Technique}: outlier handling, Hessian/reconstruction, rotation/smoothing, mixed precision, codebook/vector quantization, and data-free quantization. This axis captures the mechanism used to preserve quality.
  \item \textbf{By Scenario}: long-context LLMs, MoE LLMs, reasoning LLMs, multimodal LLMs, edge deployment, and high-throughput serving. This axis captures workload-specific constraints that cannot be inferred from bit-width alone.
\end{itemize}

These axes are complementary rather than mutually exclusive. For example, SmoothQuant~\cite{xiao2022smoothquant}\ is a PTQ method by stage, targets both weights and activations, uses transformation-based outlier smoothing, and is mainly motivated by W8A8 serving. GPTQ~\cite{frantar2022gptq}\ is PTQ by stage, weight quantization by target, Hessian-guided reconstruction by technique, and W4A16 deployment by scenario. KIVI~\cite{liu2024kivi}\ is serving-time quantization by stage, KV cache quantization by target, asymmetric low-bit scaling by technique, and long-context inference by scenario. This multi-axis view prevents a common problem in survey writing: placing each method in only one flat category and losing the reason it was proposed.

\TODO{Add a taxonomy figure with LLM Quantization at the center and four branches: By Stage, By Target, By Technique, and By Scenario.}

\section{LLM Quantization by Stage}

The stage axis organizes methods by the point at which low precision enters the model lifecycle. Earlier-stage methods can shape model parameters and training dynamics around quantization, but they require more compute and model access. Later-stage methods are cheaper and easier to apply to released models, but must compensate for quantization error with limited calibration data.

\subsection{Post-Training Quantization}
PTQ quantizes a pre-trained LLM without full retraining and is currently the most practical deployment path. It is attractive because many strong LLMs are only available as trained checkpoints and because users often lack the data or compute needed for QAT. The main PTQ challenge is to approximate the original model with a low-bit representation using only calibration samples, local reconstruction objectives, or distribution transformations.

\subsubsection{Rounding and Reconstruction}
Round-to-nearest (RTN) provides a simple baseline but ignores layer sensitivity and error propagation. GPTQ~\cite{frantar2022gptq}\ improves weight-only PTQ by using approximate second-order information to quantize columns while compensating reconstruction error. This line of work treats quantization as a local reconstruction problem: preserve each layer's output as much as possible under low-bit weights. Such methods are especially influential for W4A16 inference, where activations remain in higher precision and weight bandwidth is the primary bottleneck.

\subsubsection{Activation-Aware and Outlier-Aware PTQ}
Activation-aware methods observe that not all weights contribute equally under real inputs. AWQ~\cite{lin2023awq}\ protects a small fraction of salient weights identified through activation statistics, improving low-bit weight quantization without full retraining. LLM.int8()~\cite{dettmers2022llmint8}\ shows that emergent activation outliers can dominate quantization error and uses mixed-precision decomposition to preserve them. SmoothQuant~\cite{xiao2022smoothquant}\ instead migrates quantization difficulty from activations to weights, enabling efficient W8A8 inference. These methods are unified by a shared principle: calibration activations reveal which channels or values require special treatment.

\subsubsection{Transformation-Based PTQ}
Transformation-based PTQ reshapes weights or activations before quantization while preserving the mathematical function of the network. SmoothQuant~\cite{xiao2022smoothquant}\ uses diagonal scaling; QuIP and QuIP\#~\cite{chee2023quip,tsai2024quipsharp}\ use incoherence processing and codebooks; rotation-based methods such as QuaRot and SpinQuant~\cite{liu2024spinquant}\ reduce outlier concentration through orthogonal transformations. These methods are important because they attack the distributional cause of quantization error rather than only optimizing rounding decisions.

\subsubsection{Cross-Layer and Cross-Block PTQ}
Layer-wise PTQ can underestimate error accumulation across residual streams and attention blocks. Cross-layer and cross-block methods therefore optimize over larger units, align equivalent transformations across adjacent modules, or explicitly model propagation of quantization noise. This direction is particularly relevant for very low-bit quantization, long-context generation, and models whose capabilities depend on fragile multi-layer computations.

\subsection{Quantization-Aware Training}
QAT introduces fake quantization during training so that the model learns to tolerate low-precision representations. Compared with PTQ, QAT usually achieves better accuracy at aggressive bit-widths, but it requires training data, compute, and access to the optimization process. In LLMs, full-scale QAT is expensive, so practical QAT often appears as limited continued training, adapter-based training, or distillation-assisted calibration. QAT is most valuable when deployment requires low-bit activations, extremely low-bit weights, or hardware-native formats that are difficult to recover with PTQ alone.

\subsection{Quantized Fine-Tuning}
Quantized fine-tuning reduces the memory cost of adapting LLMs. QLoRA~\cite{dettmers2023qlora}\ freezes a 4-bit NF4 base model and trains LoRA adapters, making fine-tuning large models feasible on much smaller hardware. Follow-up methods such as QA-LoRA and LoftQ~\cite{liu2023qalora,li2023loftq}\ refine the interaction between quantization and low-rank adaptation. The key distinction from PTQ is that quantized fine-tuning does not merely compress a finished model; it changes the adaptation pipeline by deciding which tensors are frozen, which tensors receive gradients, and how quantization error affects adapter learning.

\subsection{Low-Precision Pre-Training}
Low-precision pre-training designs the model and optimizer around quantization from the beginning. FP8-LM~\cite{wu2023fp8lm}\ studies FP8 training recipes, while BitNet b1.58~\cite{ma2024bitnetb158}\ explores ternary-weight models trained from scratch. These methods are more invasive than PTQ or fine-tuning, but they may unlock fundamentally different efficiency regimes because the model never relies on high-precision weights in the same way. This stage is closely tied to hardware trends such as FP8/FP4 tensor cores and microscaling formats.

\section{LLM Quantization by Target}

The target axis organizes methods by the tensors they compress. This axis is deployment-critical because different tensors dominate different costs. Weights dominate model storage and decode-time bandwidth; activations determine whether matrix multiplication can run fully in low precision; KV cache dominates long-context memory; gradients and optimizer states dominate training memory.

\subsection{Weight Quantization}
Weight quantization is the most mature target for LLMs. Weight-only methods such as GPTQ~\cite{frantar2022gptq}, AWQ~\cite{lin2023awq}, SpQR~\cite{dettmers2023spqr}, SqueezeLLM~\cite{kim2023squeezellm}, QuIP~\cite{chee2023quip}, and AQLM~\cite{egiazarian2024aqlm} typically keep activations in FP16/BF16 while storing weights in 4-bit or lower formats. This setting is attractive because it reduces model size and memory bandwidth without requiring fully low-precision activation kernels. However, weight-only compression does not automatically guarantee speedups: dequantization overhead, packing layout, group size, and kernel support determine whether smaller weights translate into lower latency.

\subsection{Activation Quantization}
Activation quantization is necessary for end-to-end low-precision compute such as W8A8 or W4A4 inference. It is harder than weight quantization because activation distributions are input-dependent and contain channel-wise outliers. LLM.int8()~\cite{dettmers2022llmint8}\ and SmoothQuant~\cite{xiao2022smoothquant}\ are representative responses to this problem: the former isolates outlier channels, while the latter redistributes scale between activations and weights. Activation quantization is therefore inseparable from outlier handling and hardware kernel design.

\subsection{KV Cache Quantization}
During autoregressive decoding, every generated token appends keys and values to the cache. For long-context inference and high-concurrency serving, KV cache memory can exceed weight memory. KV cache quantization methods such as KIVI~\cite{liu2024kivi}, KVQuant~\cite{hooper2024kvquant}, and GEAR~\cite{zhang2024gear}\ reduce this cost by exploiting the different distributions of keys and values, asymmetric scaling, grouping, or error compensation. Unlike weight PTQ, KV cache quantization must preserve attention behavior over many future decoding steps, making long-context evaluation essential.

\subsection{Attention Quantization}
Attention quantization targets intermediate tensors in attention computation, including queries, keys, values, attention scores, and attention outputs. These tensors are sensitive because small errors in attention scores can change token routing across the sequence. Attention quantization overlaps with activation and KV cache quantization but deserves separate treatment when methods modify head-wise precision, score computation, or long-context attention mechanisms.

\subsection{Embedding Quantization}
Embedding and output projection layers can be large and semantically sensitive. Quantizing embeddings may reduce memory, but errors in token representations can propagate through the entire model. Output heads also interact with vocabulary distributions and generation quality. For this reason, many practical systems treat embeddings and final layers with higher precision or special mixed-precision rules.

\subsection{Optimizer and Gradient Quantization}
Training and fine-tuning require memory for gradients, optimizer states, and communication buffers. Quantizing these tensors differs from inference quantization because errors affect optimization dynamics rather than a single forward pass. FP8 training, 8-bit optimizers, and low-precision communication are therefore part of the broader LLM quantization landscape, especially for large-scale pre-training and distributed fine-tuning.

\section{LLM Quantization by Technique}

The technique axis explains how methods preserve model quality under low precision. This axis is often more informative than bit-width alone: two W4A16 methods can behave very differently depending on whether they use Hessian reconstruction, activation-aware scaling, rotation, sparse outlier preservation, or learned codebooks.

\subsection{Outlier Handling}
Outlier handling addresses the heavy-tailed distributions that make LLM activations and weights difficult to quantize. LLM.int8()~\cite{dettmers2022llmint8}\ explicitly separates outlier channels; SmoothQuant~\cite{xiao2022smoothquant}\ smooths activation outliers by shifting difficulty into weights; AWQ~\cite{lin2023awq}\ protects salient weights associated with high-activation channels. The shared idea is that rare but important values should not dictate the quantization scale for the entire tensor.

\subsection{Hessian and Reconstruction}
Hessian-guided and reconstruction-based methods minimize the effect of quantization on layer outputs. GPTQ~\cite{frantar2022gptq}\ uses approximate second-order information to guide weight quantization, while related block reconstruction methods optimize rounding or quantization parameters to preserve local behavior. These techniques are powerful for PTQ because they extract more information from a small calibration set than naive rounding.

\subsection{Rotation and Smoothing}
Rotation and smoothing methods transform tensor distributions before quantization. SmoothQuant~\cite{xiao2022smoothquant}\ uses scaling transformations; QuIP/QuIP\#~\cite{chee2023quip,tsai2024quipsharp}\ use incoherence processing; SpinQuant~\cite{liu2024spinquant}\ learns rotations that reduce quantization difficulty. These methods are attractive because many transformations are function-preserving or nearly function-preserving, allowing quantization to operate on better-conditioned tensors.

\subsection{Mixed Precision}
Mixed precision assigns different bit-widths to layers, channels, experts, or tensor types according to sensitivity. It is especially useful when a small subset of components dominates quantization error. In LLMs, mixed precision often appears as high-precision outlier channels, higher precision for embeddings or final layers, or separate formats for weights, activations, and KV cache. The challenge is to choose precision allocations that match both accuracy requirements and hardware efficiency.

\subsection{Codebook and Vector Quantization}
Codebook and vector quantization replace scalar uniform quantization with learned or structured codebooks. AQLM~\cite{egiazarian2024aqlm}\ uses additive quantization to represent weights with multiple codebooks, while QuIP\#~\cite{tsai2024quipsharp}\ combines incoherence with lattice codebooks. These methods are promising for 2--3 bit compression because uniform scalar quantization becomes increasingly limited at extreme bit-widths. Their practical deployment, however, depends on lookup efficiency, memory layout, and kernel support.

\subsection{Data-Free Quantization}
Data-free quantization aims to quantize models without access to representative calibration data. This setting is important when data are private, unavailable, or legally restricted. For LLMs, data-free quantization remains challenging because activation ranges, outliers, and instruction-following behavior are strongly input-dependent. A likely future direction is to combine synthetic calibration prompts, model-generated data, and conservative mixed-precision rules.

\section{LLM Quantization by Scenario}

\subsection{Long-Context LLMs}
Long-context LLMs stress quantization because KV cache memory grows linearly with sequence length. KVQuant~\cite{hooper2024kvquant}, KIVI~\cite{liu2024kivi}, KVmix~\CITE, WKVQuant~\cite{hooper2024kvquant}, ZipCache~\CITE, PQCache~\CITE, Cocktail~\CITE, and Oaken~\CITE\ address long-context memory efficiency.

\subsection{Mixture-of-Experts LLMs}
MoE LLMs require expert-aware quantization because experts differ in routing frequency, activation statistics, and sensitivity. CodeQuant~\CITE\ targets MoE outlier smoothing. DeepSeek-V3~\CITE\ demonstrates large-scale FP8 training in an MoE setting.

\subsection{Reasoning LLMs}
Reasoning models with long chain-of-thought generation may amplify quantization errors. Studies~\CITE\ suggest W8A8 and W4A16 are often safer than more aggressive low-bit settings for math and coding tasks. QeRL~\CITE\ explores quantization-enhanced reinforcement learning.

\subsection{Multimodal LLMs}
Multimodal LLMs combine visual encoders, projectors, language backbones, and sometimes audio or video modules. Q-VLM~\CITE\ studies heterogeneous precision needs in vision-language models.

\subsection{Edge Deployment}
Edge deployment emphasizes memory footprint, energy, thermal limits, and offline usability. llama.cpp/GGUF quantization types such as Q4\_K\_M, Q5\_K\_S, and Q2\_K are widely used.

\subsection{High-Throughput Serving}
High-throughput serving emphasizes tokens/s, TTFT, TPOT, batching efficiency, and GPU memory utilization. Quantization interacts with continuous batching, PagedAttention, speculative decoding, and KV cache allocation. QSpec~\CITE\ and QuantSpec~\CITE\ study quantization with speculative decoding.

\section{Systems, Hardware, and Deployment}

\subsection{Inference Frameworks}
Representative systems include llama.cpp/GGUF~\CITE, vLLM~\CITE, TensorRT-LLM~\CITE, SGLang~\CITE, ONNX Runtime, OpenVINO, Core ML, ExLlamaV2, bitsandbytes, AutoGPTQ, AutoAWQ, HQQ, and AMD Quark.

\subsection{Quantization Formats}
GGUF defines many mixed quantization types. GPTQ, AWQ, EXL2, HQQ, and bitsandbytes use different conventions. Hardware-native formats such as FP8, MXFP4, and NVFP4 are increasingly important.

\subsection{Hardware-Native Quantization}
NVIDIA Hopper supports FP8, while Blackwell introduces FP4/MXFP4/NVFP4 support. AMD MI355X supports MXFP4. Edge platforms include Apple ANE, Qualcomm Hexagon NPU, ARM CPUs, and mobile GPUs.

\subsection{Quantized Kernels}
Efficient kernels include W4A16 dequantization-fused GEMM such as Marlin, W8A8 and W4A4 integer GEMM, LUT-GEMM~\CITE, fused Hadamard-quantize-GEMM in Quartet~\CITE, QuTLASS~\CITE, Triton kernels, CUTLASS kernels, and shift-and-sum implementations.

\section{Evaluation and Benchmarking}

\subsection{Accuracy Metrics}
Perplexity on WikiText-2, C4, and PTB remains common but insufficient. LLM quantization should also be evaluated on MMLU, GSM8K, AIME, GPQA, HumanEval, LiveCodeBench, and instruction-following benchmarks.

\subsection{Long-Context Evaluation}
Long-context evaluation should include RULER, InfiniteBench, needle-in-a-haystack tests, long-document QA, and multi-turn dialogue.

\subsection{System Metrics}
System evaluation should report memory footprint, peak GPU memory, throughput, TTFT, TPOT, energy per token, batch size, context length, kernel implementation, and hardware platform.

\section{Safety, Robustness, and Reliability}

\subsection{Safety Alignment}
Quantization can degrade safety alignment even when perplexity remains acceptable. CAQ~\CITE\ introduces a contrastive alignment loss to preserve safety during PTQ.

\subsection{Adversarial Robustness}
Defensive Quantization~\CITE\ studies adversarial vulnerability under quantization. Later work reports mixed effects depending on the attack and model.

\subsection{Backdoor and Privacy Risks}
Qu-anti-zation~\CITE, Quantization Blindspots~\CITE, and EFRAP~\CITE\ study quantization-conditioned backdoors. Dynamic quantization may also leak information through shared scales or temporal scale variation~\CITE.

\section{Theoretical Understanding}

\subsection{Error Modeling}
Quantization error can be modeled as additive noise, multiplicative noise, or structured perturbation. In LLMs, errors propagate through residual streams, attention, and normalization. QEP~\CITE, LoaQ~\CITE, and LPCD~\CITE\ explicitly address error propagation.

\subsection{Scaling Laws}
Recent work studies relationships among model size, training data, group size, and precision. Topics include low-bit quantization favoring undertrained LLMs~\CITE, FP quantization scaling laws~\CITE, QAT scaling laws~\CITE, and critical data size phenomena~\CITE.

\subsection{Emergent Abilities}
Quantization can affect in-context learning, chain-of-thought reasoning, coding, mathematical problem solving, and tool use. These abilities may degrade nonlinearly with bit-width.

\section{Open Problems and Future Directions}

\begin{enumerate}[leftmargin=*,label=\textbf{\arabic*.}]
  \item \textbf{Sub-4-bit LLM quantization.} Closing the accuracy gap at 2--3 bits while preserving hardware efficiency remains challenging.
  \item \textbf{Quantization for reasoning models.} Long chain-of-thought generation may amplify quantization errors.
  \item \textbf{KV cache and long-context compression.} Future methods must reduce cache memory without harming retrieval and reasoning.
  \item \textbf{Quantization-native training.} FP8, FP4, MXFP4, ternary, and binary training require better optimization theory and hardware-software co-design.
  \item \textbf{MoE and multimodal LLM quantization.} Expert-specific and modality-specific precision allocation remains underdeveloped.
  \item \textbf{Safety-aware quantization.} PTQ and QAT objectives should explicitly preserve alignment and refusal behavior.
  \item \textbf{Standardized benchmarking.} The field needs reproducible comparisons across bit-widths, formats, kernels, hardware, tasks, and context lengths.
  \item \textbf{Interoperable formats.} Deployment would benefit from convergence among GGUF, GPTQ, AWQ, EXL2, HQQ, and hardware-native formats.
\end{enumerate}

\section{Conclusion}
This survey reorganizes LLM quantization around a taxonomy of stage, target, technique, and scenario. By focusing on LLM-specific challenges such as activation outliers, KV cache growth, long-context inference, reasoning degradation, MoE routing, multimodal heterogeneity, and safety preservation, the proposed framework separates LLM quantization from broader neural network quantization.

\section*{Acknowledgment}
\TODO{Advisor acknowledgment, funding, and collaborator acknowledgment.}

\bibliographystyle{IEEEtran}
\bibliography{references}

\balance
\end{document}